% id: 0157
% book: RPM 공통수학2 · 02 직선의 방정식
% section: 유형 UP 16 직선이 항상 지나는 점의 활용
% figure: none
% answer: ②
% answer_source: 답지
% variant_level: 0 (원본 전사)
% 이 파일은 build.mjs 가 items.json 에서 만든다. 고칠 때는 items.json 을 고친다.
\smash{\raisebox{1.5mm}{\dmkichul{RPM 공통수학2 0157번 · 27쪽}}}
직선~$y=m(x-1)+3$이 두 점~$\pt{A}(3,\,4)$, $\pt{B}(5,\,-1)$을 이은 선분~$\pt{AB}$와 만나도록 하는 실수 $m$의 값의 범위가 $\alpha\le m\le\beta$일 때, $\alpha+\beta$의 값은?
\begin{choices32}{\rule[-3ex]{0pt}{8ex}$-\dfrac{3}{2}$}{\rule[-3ex]{0pt}{8ex}$-\dfrac{1}{2}$}{\rule[-3ex]{0pt}{8ex}$\dfrac{1}{2}$}{\rule[-3ex]{0pt}{8ex}$1$}{\rule[-3ex]{0pt}{8ex}$3$}\end{choices32}
